
\begin{example}[Task Manager Engine]\label{ex:task-manager-engine}%
    Let us define a task manager engine that can create and manage tasks by
    spawning worker engines. This example demonstrates all the model
    characteristics including spawning engines, message passing, and state
    management.
    
    First, we define the message interface for the task manager engine:
    
    \begin{equation}\label{eq:task-manager-message-interface}
    \begin{aligned}
    \typedef%
    {
      \msgtype{\mathsf{taskmgr}}
    }{    &\quad \msg{CreateTask}(\N \times \N) \\
    \mid &\quad \msg{TaskStatus}(\N) \\
    \mid &\quad \msg{TaskResult}(\N \times \N) \\
    \mid &\quad \msg{WorkerReady}(\N) \\
    \mid &\quad \msg{WorkerDone}(\N \times \N)
    }\,,
    \end{aligned}
    \end{equation}
    
    where:
    \begin{itemize}
        \item $\msg{CreateTask}(n, m)$ creates a new task to compute $n + m$
        \item $\msg{TaskStatus}(id)$ queries the status of task $id$
        \item $\msg{TaskResult}(id, result)$ reports the result of task $id$
        \item $\msg{WorkerReady}(id)$ indicates a worker is ready to process tasks
        \item $\msg{WorkerDone}(id, result)$ reports a worker has completed a task
    \end{itemize}
    
    The configuration type for the task manager engine specifies its operational mode:
    
    \begin{equation}\label{eq:task-manager-configuration}
    \begin{aligned}
    \typedef{%
      \ctype{C}_{\mathsf{taskmgr}}
    }{    &\quad \ccon{Master} \\
    \mid &\quad \ccon{Worker}
    }\,.
    \end{aligned}
    \end{equation}
    
    The environment type tracks the task manager's state:
    
    \begin{equation}\label{eq:task-manager-environment}
    \begin{aligned}
    \typedef{%
      \ctype{L}_{\mathsf{taskmgr}}
    }{    &\quad \ccon{TaskManagerState}(\Map{\N}{\N} \times \Map{\N}{\Address}) \\
    \mid &\quad \ccon{WorkerState}(\N)
    }\,,
    \end{aligned}
    \end{equation}
    
    where:
    \begin{itemize}
        \item $\ccon{TaskManagerState}$ tracks task results and worker assignments
        \item $\ccon{WorkerState}$ tracks the worker's current task
    \end{itemize}
    
    Now we define the guarded actions for the task manager engine. First, the guard for creating tasks:
    
    \begin{equation}\label{eq:task-manager-create-guard}
    \begin{aligned}
    &\typing{
      g_{\mathsf{create}}
    }{
      \In{\mathsf{taskmgr}}
        \to
      \Option{\N \times \N}
    }\,. \\
    &\defequal{
      g_{\mathsf{create}}
    }{
      \lambda \langle m, c, \langle \mathsf{st}, \mathsf{workers} \rangle \rangle .
      \mathsf{case}\ m\ \mathsf{of}\ \{ \\
    &\quad \msg{CreateTask}(n, m) \Rightarrow \some{\langle n, m \rangle}; \\
    &\quad \ignore \Rightarrow \none \\
    &\}
    }\,.
    \end{aligned}
    \end{equation}
    
    The corresponding action spawns a worker engine to process the task:
    
    \begin{equation}\label{eq:task-manager-create-action}
    \begin{aligned}
    &\typing{
      a_{\mathsf{create}}
    }{
      \N \times \N \times \In{\mathsf{taskmgr}}
        \to \Out{\mathsf{taskmgr}}
    }\,. \\
    % &\defequal{
      a_{\mathsf{create}}
    % }{
      \lambda \langle n, m, \langle \ignore, \ignore, \langle \mathsf{st}, \mathsf{workers} \rangle \rangle \rangle . \\
    &\quad \chain{ \\
    &\quad \quad \spawn{ \\
    &\quad \quad \quad \configval{%
                \some{\self}%
              }{%
                \process%
              }{%
                \ccon{Worker}%
              }%
            }{%
              \localenvval{%
                \ccon{WorkerState}(0)%
              }{%
                [(\self, \ignore)]%
              }%
            }%
          }{ \\
    &\quad \quad \send{\self}{\msg{WorkerReady}(\freshid)} \\
    &\quad \quad } \\
    &\quad \}
    \,.
    \end{aligned}
    \end{equation}
    
    
    The guard for processing worker results:
    
    \begin{equation}\label{eq:task-manager-result-guard}
    \begin{aligned}
    &\typing{
      g_{\mathsf{result}}
    }{
      \In{\mathsf{taskmgr}}
        \to
      \Option{\N \times \N}
    }\,. \\
    &\defequal{
      g_{\mathsf{result}}
    }{
      \lambda \langle m, \ignore, \langle \mathsf{st}, \mathsf{workers} \rangle \rangle .
      \mathsf{case}\ m\ \mathsf{of}\ \{ \\
    &\quad \msg{WorkerDone}(id, result) \Rightarrow \some{\langle id, result \rangle}; \\
    &\quad \ignore \Rightarrow \none \\
    &\}
    }\,.
    \end{aligned}
    \end{equation}
    
    The corresponding action updates the task status:
    
    \begin{equation}\label{eq:task-manager-result-action}
    \begin{aligned}
    &\typing{
      a_{\mathsf{result}}
    }{
      \N \times \N \times \In{\mathsf{taskmgr}}
        \to \Out{\mathsf{taskmgr}}
    }\,. \\
    &\defequal{
      a_{\mathsf{result}}
    }{
      \lambda \langle id, result, \langle \ignore, \ignore, \langle \mathsf{st}, \mathsf{workers} \rangle \rangle \rangle . \\
    &\quad \update{%
            \langle
              \mathsf{st}[id \mapsto result],
              \mathsf{workers}
            \rangle
          }%
    }\,.
    \end{aligned}
    \end{equation}
    
    Finally, we define the complete behaviour of the task manager engine:
    
    \begin{equation}\label{eq:task-manager-behaviour}
    \begin{aligned}
    \defequal{
      \typing{\behaviour{\mathsf{taskmgr}}}{\behaviourtype{\mathsf{taskmgr}}}
    }{
      [ g_{\mathsf{create}} \odot a_{\mathsf{create}}
      , g_{\mathsf{result}} \odot a_{\mathsf{result}}
      ]
    }\,.
    \end{aligned}
    \end{equation}
    \end{example}